\section*{Capítulo \ref{ch:b3:solid-state} --- Química do estado sólido: bandas, defeitos e semicondutores}

\begin{solution}{exo:b3:solid-state:1}
$(E - \alpha)/|\beta| = -2\cos(k\pi/7)$: $-1.802$, $-1.247$, $-0.445$, $+0.445$, $+1.247$, $+1.802$. Eles ocupam
$3.604|\beta|$, já 90\,\% da largura de banda $4|\beta|$ da cadeia infinita.
\end{solution}

\begin{solution}{exo:b3:solid-state:2}
$2N$ ($N$ orbitais, dois elétrons cada). O sódio, com um elétron 3s por átomo, preenche metade
de sua banda s: um metal. O magnésio tem dois e preencheria exatamente a banda s,
mas as bandas 3s e 3p se sobrepõem, de modo que os níveis ocupados mais altos pertencem a uma
banda combinada parcialmente preenchida: também um metal.
\end{solution}

\begin{solution}{exo:b3:solid-state:3}
$\lambda \approx \qty{1239.8}{eV.nm}/E_g$: fosfeto de gálio \qty{549}{nm} (verde); nitreto de gálio \qty{365}{nm}
(ultravioleta próximo). Os diodos azuis usam nitreto de gálio ligado com índio, cuja
\omterm{def:b3:solid-state:band}{banda proibida} é menor.
\end{solution}

\begin{solution}{exo:b3:solid-state:4}
$\ce{CaCl2} \xrightarrow{\ce{NaCl}} \mathrm{Ca_{Na}^{\bullet} + V_{Na}' + 2\,Cl_{Cl}^{\times}}$: dois sítios de cátion e dois de ânion, como no NaCl; um Ca e
dois Cl; carga $+1 - 1 = 0$. $\ce{Y2O3} \xrightarrow{\ce{ZrO2}} \mathrm{2\,Y_{Zr}' + 3\,O_O^{\times} + V_O^{\bullet\bullet}}$: dois sítios de cátion e quatro de oxigênio; dois Y
e três O; carga $-2 + 2 = 0$.
\end{solution}

\begin{solution}{exo:b3:solid-state:5}
$T^{3/2} = 5196$ a \qty{300}{K}, $kT = \qty{0.02585}{eV}$. Silício: $\sqrt{N_cN_v} = \qty{2.42e19}{cm^{-3}}$, $\eu^{-1.125/0.05170} = \num{3.55e-10}$,
$n_i = \qty{8.6e9}{cm^{-3}}$. Germânio: $\sqrt{N_cN_v} = \qty{7.16e18}{cm^{-3}}$, $\eu^{-0.661/0.05170} = \num{2.80e-6}$, $n_i = \qty{2.0e13}{cm^{-3}}$. Razão
$\num{4.3e-4}$: a \omterm{def:b3:solid-state:band}{banda proibida} menor do germânio lhe dá mais de dois mil vezes mais
portadores.
\end{solution}

\begin{solution}{exo:b3:solid-state:6}
$n = N_D = \qty{1.0e16}{cm^{-3}}$; $p = n_i^2/n = \num{1.0e20}/\num{1.0e16} = \qty{1.0e4}{cm^{-3}}$.
\end{solution}

\begin{solution}{exo:b3:solid-state:7}
$E_F - E_i = kT\ln(n/n_i) = \qty{0.02585}{eV} \times \ln(10^6) = \qty{0.36}{eV}$.
\end{solution}

\begin{solution}{exo:b3:solid-state:8}
$kT = \qty{0.06894}{eV}$ a \qty{800}{K}; $n/N = \eu^{-2.3/0.1379} = \num{5.7e-8}$; em um mol, $\num{6.02e23} \times \num{5.7e-8} = \num{3.4e16}$
pares.
\end{solution}

\begin{solution}{exo:b3:solid-state:9}
A massa de uma cela é $\rho a^3 = \qty{5.70}{g.cm^{-3}} \times (\qty{4.300e-8}{cm})^3 = \qty{4.532e-22}{g}$, isto é,
$\qty{4.532e-22}{g} \times N_A = \qty{272.9}{g}$ por mol de celas, ou \qty{68.23}{g} por oxigênio. Então
$(1 - x) \times 55.8 + 16.0 = 68.23$ dá $1 - x = 0.936$: $x = 0.064$, $\mathrm{Fe_{0.936}O}$.
\end{solution}

\begin{solution}{exo:b3:solid-state:10}
$\ln(\sigma T)$: $-3.51$, $-1.30$, $0.36$ para $1/T = 3.333$, $2.857$, $\qty{2.500e-3}{K^{-1}}$; inclinação entre os pontos extremos
$(0.36 + 3.51)/(-\qty{8.33e-4}{K^{-1}}) = -\qty{4.65e3}{K}$ (o ponto do meio fica sobre a reta), logo
$E_a = \qty{4.65e3}{K} \times \qty{8.617e-5}{eV.K^{-1}} = \qty{0.40}{eV}$.
\end{solution}

\begin{solution}{exo:b3:solid-state:11}
$\tfrac12\ce{O2} \rightleftharpoons \mathrm{O_O^{\times} + V_M' + h^{\bullet}}$, $K = [\mathrm{V_M'}][\mathrm h^\bullet]/p(\ce{O2})^{1/2}$. Balanço de carga $[\mathrm h^\bullet] = [\mathrm{V_M'}]$, logo
$[\mathrm h^\bullet]^2 = K\,p(\ce{O2})^{1/2}$ e $[\mathrm h^\bullet] \propto p(\ce{O2})^{1/4}$; a condutividade acompanha os \omterm{def:b3:solid-state:carriers}{buracos}.
\end{solution}

\begin{solution}{exo:b3:solid-state:12}
$n/N = \sqrt{N_i/N}\,\eu^{-\Delta H_F/2kT} = \sqrt 2\,\eu^{-1.4/(2 \times 0.05170)} = 1.414 \times \num{1.32e-6} = \num{1.9e-6}$.
\end{solution}

\begin{solution}{pb:b3:solid-state:1}
\textbf{1.} $\ce{Y2O3} \xrightarrow{\ce{ZrO2}} \mathrm{2\,Y_{Zr}' + 3\,O_O^{\times} + V_O^{\bullet\bullet}}$.

\textbf{2.} Sítios: dois sítios de cátion e quatro sítios de oxigênio (três ocupados, um vazio),
a razão 1\,:\,2 do $\ce{ZrO2}$. Massa: 2 Y e 3 O de cada lado. Carga: $2 \times (-1) + 2 = 0$.

\textbf{3.} Uma lacuna para dois íons ítrio: meia por ítrio.

\textbf{4.} Para $(\ce{ZrO2})_{0.92}(\ce{Y2O3})_{0.08}$ há 0.92 Zr e 0.16 Y, 1.08 cátion: $x = 0.16/1.08 = 0.148$,
$\mathrm{Zr_{0.852}Y_{0.148}O_{1.926}}$.

\textbf{5.} $(x/2)/2 = 0.0370$: 3.7\,\% dos sítios de oxigênio estão vazios.

\textbf{6.} $8 \times 0.0370 = 0.296$ lacuna por cela.

\textbf{7.} $a^3 = (\qty{5.14e-8}{cm})^3 = \qty{1.358e-22}{cm^3}$: $0.296/\qty{1.358e-22}{cm^3} = \qty{2.2e21}{cm^{-3}}$.

\textbf{8.} Cada íon óxido precisa saltar uma barreira de \qty{1.0}{eV} para uma lacuna
vizinha; a fração de tentativas bem-sucedidas, $\eu^{-E_a/kT}$, cresce muito depressa com $T$.

\textbf{9.} A \qty{1073.15}{K}, $kT = \qty{0.09248}{eV}$: $A = \sigma T\eu^{E_a/kT} = 0.030 \times 1073.15 \times \eu^{10.81} = \qty{1.6e6}{S.K.cm^{-1}}$.

\textbf{10.} A \qty{973.15}{K}: $\sigma = (A/T)\eu^{-11.92} = \qty{0.011}{S.cm^{-1}}$.

\textbf{11.} A \qty{573.15}{K}: $\sigma = (A/T)\eu^{-20.25} = \qty{4.5e-6}{S.cm^{-1}}$.

\textbf{12.} $R = L/(\sigma S) = \qty{0.10}{cm}/(\sigma \times \qty{0.20}{cm^2})$: \qty{46}{\ohm} a \qty{700}{\celsius}, $\qty{1.1e5}{\ohm}$ a \qty{300}{\celsius}.

\textbf{13.} $R < \qty{1}{k\ohm}$ exige $\sigma > \qty{5.0e-4}{S.cm^{-1}}$; resolver $(A/T)\eu^{-E_a/kT} = \num{5.0e-4}$ (por tentativas)
dá $T \approx \qty{760}{K}$, cerca de \qty{490}{\celsius}.

\textbf{14.} O gás de escape está frio após a partida e em baixa carga; o aquecedor leva o
eletrólito à sua temperatura de operação em segundos, de modo que o controle
funciona desde a partida, quando a maior parte dos poluentes é emitida.

\textbf{15.} $\ce{O2 + 4e- -> 2O^2-}$; em \omterm{def:b3:solid-state:kroger-vink}{notação de Kröger--Vink}, $\ce{O2} + \mathrm{2\,V_O^{\bullet\bullet} + 4\,e'} \longrightarrow \mathrm{2\,O_O^{\times}}$.

\textbf{16.} O oxigênio é reduzido do lado do ar, e os íons óxido são reoxidados
a oxigênio do lado do gás de escape: a célula transfere $\ce{O2}$ do ar para o gás de escape, com
$\Delta_rG = RT\ln(p_{\mathrm{exh}}/p_{\mathrm{air}})$ por mol de $\ce{O2}$ e quatro elétrons: $E = -\Delta_rG/4F = (RT/4F)\ln(p_{\mathrm{air}}/p_{\mathrm{exh}})$.

\textbf{17.} $RT/4F = 8.3145 \times 973.15/(4 \times 96485) = \qty{0.02096}{V}$.

\textbf{18.} Zero: as duas pressões são iguais.

\textbf{19.} $(RT/4F)\ln 10 = \qty{48.3}{mV}$ por década.

\textbf{20.} $\qty{0.02096}{V} \times \ln(0.2095/0.010) = \qty{0.02096}{V} \times 3.04 = \qty{0.064}{V}$.

\textbf{21.} $\qty{0.02096}{V} \times \ln(0.2095/\num{1.0e-20}) = \qty{0.02096}{V} \times 44.49 = \qty{0.93}{V}$.

\textbf{22.} O gás de escape rico contém monóxido de carbono, hidrogênio e hidrocarbonetos
não queimados; no eletrodo de platina, o pouco oxigênio que resta reage com eles,
e $p(\ce{O2})$ é fixada pelos equilíbrios $\ce{2CO + O2 <=> 2CO2}$ e $\ce{2H2 + O2 <=> 2H2O}$, que deixam um valor
ínfimo a \qty{700}{\celsius}.

\textbf{23.} $p = 0.2095\,\eu^{-0.45/0.02096} = \qty{1.0e-10}{bar}$.

\textbf{24.} Em torno da razão estequiométrica, o gás de escape passa de um leve excesso
de oxigênio a um leve excesso de CO e de hidrogênio: $p(\ce{O2})$ cai cerca de dezoito
potências de dez para uma pequena variação de combustível, e a tensão salta de menos de
\qty{0.1}{V} para cerca de \qty{0.9}{V}. O controle do motor adiciona combustível quando a tensão está baixa e
o retira quando está alta, mantendo a mistura no ponto estequiométrico em que o
catalisador de três vias converte ao mesmo tempo CO, hidrocarbonetos e óxidos de nitrogênio.

\textbf{25.} Em gás de escape rico a \qty{700}{\celsius}, o sensor dá cerca de \qty{0.93}{V}.
\end{solution}
